%intuitive-mlengineer-robotics-computervision
\documentclass[letterpaper,10pt]{article}

\usepackage[left=0.35in, right=0.35in, top=0.35in, bottom=0.25in]{geometry}
\usepackage[inline]{enumitem}
\usepackage{hyperref}
\usepackage{titlesec}
\usepackage{xcolor}
\usepackage{environ}
\usepackage{etoolbox}
\usepackage{fontawesome5}
%\usepackage{hrefs}
\usepackage{hyperref}

% Custom colors
\definecolor{darkblue}{RGB}{0,0,139}
\hypersetup{
    colorlinks=true,
    linkcolor=darkblue,
    filecolor=darkblue,      
    urlcolor=darkblue,
}

% Section formatting - OPTIMIZED SPACING
\titleformat{\section}{\large\bfseries\uppercase}{}{0em}{}[\titlerule]
\titlespacing*{\section}{0pt}{8pt}{4pt}

% Custom commands
\newcommand{\resumeItem}[1]{\item\small{#1}}
\newcommand{\resumeSubheading}[4]{
  \vspace{-1pt}\item
    \begin{tabular*}{0.97\textwidth}[t]{l@{\extracolsep{\fill}}r}
      \textbf{#1} & #2 \\
      \textit{#3} & #4 \\
    \end{tabular*}\vspace{-3pt}
}
\newcommand{\resumeSubSubheading}[2]{
    \item
    \begin{tabular*}{0.97\textwidth}{l@{\extracolsep{\fill}}r}
      \textit{\small#1} & \small #2 \\
    \end{tabular*}\vspace{-3pt}
}
\newcommand{\resumeProjectHeading}[2]{
    \item
    \begin{tabular*}{0.97\textwidth}{l@{\extracolsep{\fill}}r}
      \textbf{#1} & #2 \\
    \end{tabular*}\vspace{-3pt}
}
\newcommand{\resumeSubItem}[1]{\resumeItem{#1}\vspace{2pt}}
\renewcommand\labelitemii{$\vcenter{\hbox{\tiny$\bullet$}}$}

% OPTIMIZED BULLET SPACING
\newcommand{\resumeSubHeadingListStart}{\begin{itemize}[leftmargin=0.15in, label={}, itemsep=1pt, parsep=0pt]}
\newcommand{\resumeSubHeadingListEnd}{\end{itemize}}
\newcommand{\resumeItemListStart}{\begin{itemize}[leftmargin=0.2in, itemsep=1pt, parsep=0pt]}
\newcommand{\resumeItemListStartEdu}{\begin{itemize}[leftmargin=0.35in, itemsep=1pt, parsep=0pt]} % <--- NEW LINE HERE
\newcommand{\resumeItemListEnd}{\end{itemize}\vspace{-2pt}}

\begin{document}

% Removes page numbers
\pagestyle{empty}

%----------HEADING----------
\begin{center}
    \textsc{\textbf{\LARGE Tushar Nayak}} \\ \vspace{5pt}
    \small (412) 478-5677 $|$ \href{mailto:tusharn@andrew.cmu.edu}{\underline{tusharn@andrew.cmu.edu}} $|$ 
    \href{https://linkedin.com/in/nayaktushar}{\underline{linkedin.com/in/nayaktushar}} $|$
    \href{https://tushar-nayak.github.io/}{\underline{tushar-nayak.github.io}}  
\end{center}

\vspace{-2pt}
\noindent\small
\noindent\small
Computer Vision and Deep Learning Researcher/Engineer focused on \textbf{multi-modal medical imaging}, 2D--3D registration and reconstruction, geometric \& physics-informed deep learning. Building real-time medical systems that recover structure from sparse visual observations.

\vspace{-2pt}

%-----------EDUCATION-----------
\section{Education}
  \resumeSubHeadingListStart
    \resumeSubheading
      {Carnegie Mellon University}{Pittsburgh, PA}
      {Master of Science - Biomedical Engineering (Research)}{May 2026}
      \vspace{-0.5pt}
      \resumeItemListStartEdu % <--- CHANGED HERE
        \resumeItem{\textsc{Courses:} Computer Vision, Learning for 3D Vision, Image-Based Computational Modeling \& Analysis, Medical Image Analysis.}
        \resumeItem{\textsc{TA:} Computer Vision, Applied Deep Learning, ML in Biomedical Research, Fundamentals of Computational Bioengineering.}
        \resumeItem{\textsc{Awards and Leadership:} CMU-BME Department Head's Fellowship, CMU-BME Department Ambassador.}
      \resumeItemListEnd
      
    \resumeSubheading
      {Manipal Institute of Technology}{Manipal, India}
      {Bachelor of Technology - Biomedical Engineering, Minor in Data Science}{May 2023}
      \vspace{-0.5pt}
      \resumeItemListStartEdu % <--- CHANGED HERE (commented out)
        \resumeItem{\textsc{Awards:} Best Paper - 2nd International Conference on Artificial Intelligence, Computational Electronics \& Communication}
        \resumeItem{\textsc{Leadership:} IEEE-EMBS Student Chapter Chair \& Research Head; IEEE-RAS Student Chapter Exec. Board \& Webmaster.}
      \resumeItemListEnd
  \resumeSubHeadingListEnd

%-----------SKILLS-----------
% \section{Skills}
% \begin{itemize}[leftmargin=0.15in, label={}, itemsep=1pt]
% \small\item{
% \textbf{Deep Learning:} GNNs, Neural ODEs, PINNs, Attention Mechanisms \& Transformers, U-Nets, VLMs, Occupancy Networks. \\
% \textbf{3D Vision:} Reconstruction, Neural Rendering, Deformable Registration, Implicit Representations, Differential Rendering. \\
% \textbf{Computer Vision:} Detection, Segmentation, Multi-view Geometry, Camera Geometry, Optical Flow, Structure from Motion. \\
% \textbf{Frameworks:} PyTorch, PyTorch3D, PyGeometric, TorchDiffEq, TensorFlow, Keras, CUDA, OpenCV, ITK, VMTK, MONAI, VTK, Git, Docker, Linux. \\
% \textbf{Languages:} Python, \textsc{MATLAB}, R, Julia, Bash, \LaTeX.
% }
% \end{itemize}
\section{Skills}
\begin{itemize}[leftmargin=0.15in, label={}, itemsep=1pt]
\small\item{
\textbf{Medical Imaging Toolkits:} MONAI, ITK / SimpleITK / ITK-Snap, VTK, VMTK, OpenCV, PyTorch3D, FSL. \\
\textbf{Deep Learning:} PyTorch, TensorFlow, GNNs, Neural ODEs, PINNs, Attention Mechanisms, Transformers, U-Nets, Occupancy Nets. \\
\textbf{Computer Vision:} Detection, Segmentation, 3D Reconstruction, Multi-view Geometry, Camera Geometry, Deformable Registration, Implicit Representation, Gaussian Splatting, Neural Rendering, Differentiable Rendering. \\
\textbf{Software:} Python, MATLAB, Bash, Git, Docker, Linux, CUDA, \LaTeX, R.}
\end{itemize}



%-----------EXPERIENCE-----------
\section{Research}
  \resumeSubHeadingListStart

    \resumeSubheading
      {Computational Engineering and Robotics Lab, Carnegie Mellon University}{Pittsburgh, PA}
      {Graduate Researcher, Real-time 3D Vascular Deformation Estimation \href{https://tushar-nayak.github.io/assets/pdf/tusharn-poster.pdf}{\textbf{} \faChalkboardTeacher}}{Aug 2024 -- present}
      \resumeItemListStart
        \resumeItem{Architected {\textsc{FluoroGRID}}, a \textbf{multi-modal quad-encoder graph neural network} fusing live 2D fluoroscopy feeds with 3D pre-operative volumes to \textbf{map 3D deformations from 2D input in real time} to reduce continuous operative radiation.}
        \resumeItem{Engineered a geometrically regularized loss function leveraging discrete differential geometry to enforce structural limits, achieving \textbf{sub-millimeter precision} and sub 0.01\% geometric plausibility penalties, at real-time inference speeds (\textbf{\textgreater~30+fps}, 28x faster).}
        \resumeItem{Built a physics-grounded synthetic data pipeline \textbf{simulating viscoelastic tool-tissue interactions} using {\texttt{vmtk}} and \textbf{fluoroscopy ray-tracing} using {\texttt{gvxr}}, generating an expansive dataset of \textbf{370,000+ 3D--2D training data pairs}.}
      \resumeItemListEnd

    \resumeSubheading
      {The $\forall$ Lab and Image Science Lab, Carnegie Mellon University}{Pittsburgh, PA}
      {Research Assistant, Longitudinal Glioblastoma Evolution Visual Forecasting {\href{https://tushar-nayak.github.io/glioblastoma-evolution/}{\small \texttt{\faLink}}}}{Feb 2025 -- Aug 2025}
      \resumeItemListStart
        \resumeItem{Developed a \textbf{Neural ODE}-based forecasting pipeline to visually predict longitudinal glioblastoma growth using multimodal MRI.}
        \resumeItem{Built an \textbf{Attention U-Net \& temporal embedding model} to simulate the evolution of latent tumor in continuous time.}
        \resumeItem{Presented findings at the \textbf{CMU BME Forum 2025} \href{https://www.youtube.com/watch?v=Rw9KHNOaAL0}{\faYoutube}, currently modeling tumor evolution as an explicit dynamics problem.}
      \resumeItemListEnd

    \resumeSubheading
      {Biomedical Informatics Computing Lab, Manipal Institute of Technology}{Manipal, India}
      {Undergraduate Researcher, Multi-Stage Multi-Modal Oral Cancer Detection}{Jan 2022 -- May 2023}
      \resumeItemListStart
        \resumeItem{Ensembled \textbf{three transfer-learning models with self-attention} and CAM across lesion photo, OCT, Histo modalities.}
        %\resumeItem{Achieved \textbf{95\% accuracy} across lesion photos, OCT, and histological modalities, with class activation mapping for explainability.}
        \resumeItem{Published \& presented research findings on this pipeline, alongside diagnostic models for Lung Cancer, Mpox, \& Dengue* \href{https://scholar.google.com/citations?user=9xUX7NoAAAAJ&hl=en}{\faGraduationCap}.}
      \resumeItemListEnd
      
  \resumeSubHeadingListEnd
    %-----------PROJECTS-----------
\section{Projects}
  \resumeSubHeadingListStart

\resumeProjectHeading
    {\textit{Sparse 2D Echocardiograph to 3D Cardiac Volume Reconstruction} {\href{https://tushar-nayak.github.io/cardiac-volume-reconstruction/}{\small{\faLink}}} {\textnormal{(CMU)}} }{Aug 2025 - present}
    \resumeItemListStart
        \resumeItem{Developed \textbf{coordinate-based neural implicit representation} model for cardiac volume reconstruction with surface priors from only three sparse 2D echocardiography slices \textbf{achieving 3D Dice 0.864 and 3D IoU 0.765} on the \textsc{Mitea CAP} dataset.}
        \resumeItem{Implemented \textbf{differentiable multi-view projection with learnable pose optimization} for joint slice alignment \& geometry recovery, Fourier positional encoding for high-frequency boundaries in sparse supervision, and curriculum \& meta learning approaches.}
        \resumeItem{Currently exploring a \textbf{stabilized Gaussian occupancy field} pipeline, combining differentiable rasterization and pose optimization.}
        %\resumeItem{Applied volumetric BCE point sampling to supervise unobserved 3D regions and geometric regularization (surface priors) to enforce anatomically plausible reconstructions, .}
    \resumeItemListEnd

  %   \resumeProjectHeading{\textit{3D Lung CT Segmentation and Surface Export {\href{https://github.com/tushar-nayak/lungvolseg}{\small \faLink}}} (CMU)}{Mar 2026 -
  % present}
  % \resumeItemListStart
  % \resumeItem{Building a \textbf{lung CT segmentation} pipeline using \textbf{\textsc{SimpleITK}} preprocessing, \textbf{MONAI 3D U-Net} and \textbf{VTK mesh extraction}. }
  % \resumeItem{Achieved \textbf{0.9243 mean Dice} and \textbf{10.59 mean HD95} on a 20-case subset of the \textsc{Zenodo COVID19 CT} dataset.}
  % \resumeItemListEnd

\resumeProjectHeading{\textit{3D Lung CT Segmentation {\href{https://tushar-nayak.github.io/lungvolseg/}{\small \faLink}}} {\textnormal{(CMU)}}}{Mar 2026 -- Apr 2026}
\resumeItemListStart
  \resumeItem{Built a \textbf{CT-to-3D anatomy pipeline} using \textbf{SimpleITK} preprocessing, \textbf{MONAI 3D U-Net} segmentation, and \textbf{VTK surface extraction} to convert volumetric lung CT into navigable 3D meshes for planned \textbf{centerline based airway modelling}.}
  \resumeItem{Achieved \textbf{0.9243 mean Dice} and \textbf{10.59 mean HD95} on the \textsc{Zenodo COVID19 CT Lung \& Infection Segmentation} dataset.}
\resumeItemListEnd

    \resumeProjectHeading{\textit{Deformable 3D Reconstruction of Live Endoscopic Surgical Scenes {\href{https://github.com/tushar-nayak/endo-splat}{\small \faLink}} {\textnormal{(CMU)}}}}{Apr 2026 - present}
    \resumeItemListStart
    \resumeItem{Prototyping a \textbf{deformable 3D Gaussian Splatting pipeline with VLM semantic embeddings} for endoscopic scenes.}
    \resumeItem{Preliminary results show a +8.5dB PSNR, 75\% depth error reduction, and initial zero-shot text querying on \textsc{EndoScapes} dataset.}
    \resumeItemListEnd
        \resumeProjectHeading
    {\textit{Dual View Mammogram-based Breast Cancer Prediction} {\href{https://github.com/tushar-nayak/grading-cbisddsm/}{\small \faLink}} {\textnormal{(CMU)}}}{Jan 2026 - May 2026}
    \resumeItemListStart
        \resumeItem{Developed a \textbf{spatial alignment pipeline} to register MLO mammograms to CC view using breast-mask \textbf{centroid matching and affine translation}, for better dual-view feature correspondence by \textbf{reducing cross-view centroid distance by 99.91\%}.}
        \resumeItem{Built a \textbf{cross-view attention classifier} that fused aligned mammograms for binary prediction at \textbf{80\%} on \textsc{cbis-ddsm} dataset.}
    \resumeItemListEnd

    
    \resumeProjectHeading
    {\textit{Neural PDE-based Anisotropic Diffusion for MRI Denoising} {\href{https://tushar-nayak.github.io/neural-anisotropic-diffusion/}{\small \faLink}} {\textnormal{(CMU)}}}{Feb 2026- Mar 2026}
    \resumeItemListStart
        %\resumeItem{Unrolled the Perona-Malik denoising PDE into a neural network with \textbf{learned spatially-adaptive conduction weights}.}
        %\resumeItem{Achieved \textbf{24.85 dB PSNR} (+4.46 dB vs. baselines) while \textbf{reducing Edge MSE by 47\%} on \textsc{Br35h} dataset.}
    \resumeItem{Unrolled Perona-Malik PDE with \textbf{learned conduction weights}, achieving \textbf{+23\% PSNR} and \textbf{47\% lower Edge MSE} on \textsc{Br35h}.}
    \resumeItemListEnd
    
    
    %resumeProjectHeading{\textit{Perception and Visual Learning {\href{https://github.com/openhorizonrobotics/perception}{\small \faLink}} {\textnormal{(Education Team, Open Horizon Robotics)}}}}{Mar 2026 - present} 
    %\resumeItemListStart
    %    \resumeItem{Authoring a \textbf{comprehensive open-source computer vision course} covering Classical 2D Vision to Localization \& Mapping.}
    %\resumeItemListEnd

    
    \resumeSubHeadingListEnd

\end{document}